\begin{abstract}
Sharpness-aware training and second-order post-training pruning use different local descriptions of a network. The former allocates perturbations according to task gradients, whereas the latter reconstructs layer outputs after constrained weight deletions. We investigate whether conditional reconstructability can connect these stages. Historical Compensability-Guided SAM (HC-SAM) estimates damped activation inverse-Gram diagonals during dense training, accumulates their globally centered log scales, and uses the resulting history to redistribute a fixed global $L_2$ perturbation budget. The trained dense model is subsequently compressed by native SparseGPT, without changing its dynamic mask selection or compensation and without recovery training. We characterize the conditional single-deletion interpretation, geometric averaging, and budget competition underlying the method. These properties do not guarantee improved accuracy or solve the full joint $N{:}M$ deletion problem. The framework separates the proposed training signal from the downstream compression algorithm, enabling its effect to be tested without changing native mask selection or compensation. The central empirical question is whether this training-side signal improves compatibility with an unchanged post-training compressor.
\end{abstract}
